\documentclass[11pt,a4paper]{article}
\usepackage[T1]{fontenc}
\usepackage{lmodern,amsmath,amssymb,amsthm}
\usepackage[a4paper,margin=26mm,headheight=14pt]{geometry}
\usepackage{microtype,booktabs,tabularx,array,graphicx,xcolor,fancyhdr}
\usepackage[colorlinks=true,linkcolor=blue!45!black,citecolor=blue!45!black,urlcolor=blue!45!black]{hyperref}
\hypersetup{pdftitle={Symbolic Resonance Geometry and Quantum Ring Spectra},
pdfauthor={Hilmir Frimann Halldorsson},
pdfsubject={Review draft v0.1, 13 October 2026: exact discrete-to-continuum spectral correspondence}}
\pagestyle{fancy}\fancyhf{}
\fancyhead[L]{\small SRG and quantum ring spectra}
\fancyhead[R]{\small Review draft v0.1}
\fancyfoot[C]{\thepage}
\setlength{\parskip}{4pt}
\setlength{\parindent}{1.1em}
\setlength{\emergencystretch}{1.5em}
\clubpenalty=10000\widowpenalty=10000
\renewcommand{\arraystretch}{1.15}
\newtheorem{theorem}{Theorem}[section]
\newtheorem{proposition}[theorem]{Proposition}
\newtheorem{corollary}[theorem]{Corollary}
\theoremstyle{definition}\newtheorem{definition}[theorem]{Definition}
\theoremstyle{remark}\newtheorem{remark}[theorem]{Remark}
\newcommand{\C}{\mathbb C}\newcommand{\R}{\mathbb R}\newcommand{\Z}{\mathbb Z}
\newcommand{\Tr}{\operatorname{Tr}}
\newcommand{\ee}{\mathrm e}\newcommand{\ii}{\mathrm i}
\newcommand{\phig}{\varphi}\newcommand{\AN}{A_{N,L}}
\newcommand{\ip}[2]{\langle #1,#2\rangle_h}
\numberwithin{equation}{section}
\title{\vspace{-12mm}\textbf{Symbolic Resonance Geometry\\and Quantum Ring Spectra}\\[5pt]
\large An Exact Discrete-to-Continuum Spectral Correspondence}
\author{Hilmir Fr\'imann Halld\'orsson}
\date{13 October 2026\\[3pt]\small Complete review edition v0.1; submitted to GPT for scientific review}
\begin{document}
\maketitle
\begin{abstract}
The periodic-loop diagnostic developed in the October 2026 reconstruction of
Symbolic Resonance Geometry (SRG) admits a direct spectral correspondence with
the kinetic Hamiltonian of a free quantum particle on a circle. For circumference
\(L>0\) and \(N\geq3\) equally spaced sites, the positive nearest-neighbor
Laplacian has eigenvalues
\(4N^2L^{-2}\sin^2(\pi k/N)\).
We derive this spectrum and prove that its nonzero inverse-power trace at
exponent \(3/2\) converges to \(L^3\zeta(3)/(4\pi^3)\).
The proof counts both Fourier directions, handles the even-\(N\) midpoint,
and supplies a finite-size error bound. Adopting the conventional kinetic
scale \(\hbar^2/(2m)\) gives an exact finite-dimensional quantum Hamiltonian
and the same limiting spectral-zeta trace as the continuum ring.
We also establish a quantitative correspondence of the low-frequency modes
and their standard unitary evolution. At unit circumference, an explicitly
calibrated multiple of the geometric trace equals SRG's historically selected
reference \(\gamma_*=\zeta(3)/(\pi e\varphi)\).
This relates an SRG-specific diagnostic to established spectral and quantum
mathematics without changing the historical recursion. Positive inverse traces,
negative finite parts and physical Hamiltonian interpretation are distinguished
by their definitions. The contribution is a source-traced correspondence with
self-contained proofs, rather than a new origin of quantum mechanics or a
first-principles selection of the historical calibration.
\end{abstract}
\noindent\textbf{Keywords:} Symbolic Resonance Geometry; periodic Laplacian;
quantum ring; discrete Fourier modes; spectral zeta; continuum limit.

\section{A geometric route to the quantum ring}\label{sec:intro}
A circle permits a wave to close after an integer number of turns. Sampling
that circle at equally spaced sites preserves periodicity but changes the
response to short wavelengths. These two facts provide the connection developed
here: the same Fourier labels diagonalize a periodic nearest-neighbor
difference operator and the free-particle kinetic operator on a continuous
ring. Their spectra agree in the refinement limit, with a computable
difference at every finite resolution.

In the SRG research program, geometry, local response and recursive
composition have been studied through several historically distinct
definitions. The completed \emph{SRG Origins} reconstruction introduced a
normalized periodic-loop diagnostic as one precise way to examine the
appearance of \(\zeta(3)\) in the source record
\cite[section 5]{origins}. That diagnostic supplies this paper's geometric
starting point. It is a spatial operator on a fixed loop. No assertion about
a unique historical temporal recursion is needed to analyze it.

The key result is
\begin{equation}
 \lim_{N\to\infty}\Tr{}'(\AN^{-3/2})
   =\Tr{}'\!\left[\left(-\frac{d^2}{dx^2}\right)^{-3/2}\right]
   =\frac{L^3}{4\pi^3}\zeta(3).
 \label{eq:headline}
\end{equation}
Here the prime means that the constant, zero-eigenvalue mode is removed and
the continuum operator has periodic boundary conditions on \([0,L]\).
The inverse power is defined on the orthogonal complement of constants.
The equality concerns complete spectral traces, not only the convergence of
one eigenvalue at a time.

Why does the zeta value have argument three? On a continuous circle the
nonzero Laplace eigenvalues grow as \(k^2\). Taking their inverse power
\(3/2\) yields a weight proportional to \(k^{-3}\). Each positive integer
frequency has two independent directions, so both signs of \(k\) must be
counted. The circumference fixes the remaining scale. Thus the number,
multiplicity and dimensions of the result can all be read from the geometry
once the observable has been declared.

The physical comparison is equally explicit. Standard quantum mechanics
uses \(H=(\hbar^2/2m)(-d^2/dx^2)\) on the periodic circle
\cite{zwiebach}. Multiplying the discrete diagnostic by this same positive
constant produces a conventional discrete kinetic Hamiltonian. Its spectral
trace is then related exactly to the geometric one, and its low-frequency
energies converge to the ring energies. Section~\ref{sec:quantum} develops
this correspondence, including the time law that follows when the
Schr\"odinger equation is adopted.

These ingredients have established mathematical precedents. Friedli and
Karlsson give spectral zeta functions for cyclic graphs, their relation to
continuous tori and the continuum circle zeta function
\cite[pp.~2--3, Eqs.~(1)--(2)]{friedli}.
Karlsson also presents the circle through its heat kernel and Mellin
transform \cite[section 2.1]{karlsson}.
We give direct proofs in the present normalization, then connect them to the
specific SRG reference and its documented history.
The Fourier spectrum, ring Hamiltonian and spectral-zeta framework are not
claimed as SRG inventions.

The paper proceeds from the finite loop and its quadratic form to the trace
limit, then to the quantum operator and calibration. A final comparison of
spectral observables explains why the positive number in
\eqref{eq:headline} and the negative finite part \(-1/12\) are compatible
without being interchangeable.

\section{The periodic geometric diagnostic}\label{sec:discrete}
\subsection{Spacing, state space and sign}
Fix a circumference \(L>0\), an integer \(N\geq3\), and spacing \(h=L/N\).
The sites are \(x_j=jh\), with \(j\) interpreted modulo \(N\).
A state is a vector \(u=(u_0,\ldots,u_{N-1})\in\C^N\), and the natural
quadrature inner product is
\begin{equation}
 \ip{u}{v}=h\sum_{j=0}^{N-1}\overline{u_j}v_j.
 \label{eq:inner}
\end{equation}
The real state space \(\R^N\) is an invariant subspace of the operator below.
Complex coordinates provide the most economical Fourier basis; their use
at this stage is algebraic.

\begin{definition}[Positive periodic Laplacian]\label{def:lap}
The adopted diagnostic is
\begin{equation}
 (\AN u)_j=\frac{2u_j-u_{j-1}-u_{j+1}}{h^2}
          =\frac{N^2}{L^2}(2u_j-u_{j-1}-u_{j+1}).
 \label{eq:lap}
\end{equation}
This is the positive sign convention, approximating \(-d^2/dx^2\).
\end{definition}
The spacing is arc length along the circle, not the straight chord between
neighboring sites in a drawing. In particular \(L\) is the circumference,
not the radius. The geometry specifies a closed one-dimensional chain;
no embedding-dependent curvature potential is included.

\begin{proposition}[Positivity and the zero mode]\label{prop:positive}
The operator \(\AN\) is self-adjoint for \eqref{eq:inner}, and
\begin{equation}
 \ip{u}{\AN u}=\frac1h\sum_{j=0}^{N-1}|u_{j+1}-u_j|^2\geq0.
 \label{eq:form}
\end{equation}
Its kernel consists exactly of constant vectors.
\end{proposition}
\begin{proof}
Expanding the right side gives
\(h^{-1}\sum_j(2|u_j|^2-\overline{u_{j+1}}u_j-
\overline{u_j}u_{j+1})\). Shifting periodic indices identifies this
with the left side. The same index shift proves
\(\ip{u}{\AN v}=\ip{\AN u}{v}\), or equivalently the matrix is real
symmetric. The form vanishes exactly when every adjacent difference
vanishes. The cycle is connected, so all entries are then equal.
\end{proof}
The restriction \(N\geq3\) ensures two distinct neighbors at each site.
For \(N=3\) each site is adjacent to the other two and the formula is
unchanged. With two sites the two offsets coincide and the same written
stencil would count that neighbor twice, describing a different graph
convention. We do not use that extension.

The inverse in \eqref{eq:headline} is therefore taken on
\begin{equation}
 \mathcal H_N^0=\left\{u\in\C^N:\sum_j u_j=0\right\}.
 \label{eq:meanfree}
\end{equation}
Removing the zero mode is essential to a finite inverse trace. It is a
definition of the observable, not a deletion of a quantum state from the
full Hamiltonian.

\subsection{Fourier modes and their multiplicities}
\begin{proposition}[Complete finite spectrum]\label{prop:fourier}
The vectors
\begin{equation}
 v_{N,k}(j)=L^{-1/2}\exp(2\pi\ii kj/N),\qquad 0\leq k<N,
 \label{eq:fourier}
\end{equation}
form an orthonormal basis for \eqref{eq:inner}. Their eigenvalues are
\begin{equation}
 \lambda_{N,k}(L)=\frac{4N^2}{L^2}\sin^2\frac{\pi k}{N}.
 \label{eq:spectrum}
\end{equation}
The zero mode is simple. Nonzero eigenvalues are paired at \(k,N-k\);
when \(N\) is even, \(k=N/2\) is an unpaired eigenmode with eigenvalue
\(4N^2/L^2\).
\end{proposition}
\begin{proof}
For integers \(k,\ell\) in the given range,
\[
 \ip{v_{N,k}}{v_{N,\ell}}
 =\frac1N\sum_{j=0}^{N-1}\exp(2\pi\ii(\ell-k)j/N).
\]
This is one if \(k=\ell\); otherwise the geometric sum has numerator
\(1-\exp(2\pi\ii(\ell-k))=0\) and a nonzero denominator. There are \(N\)
orthonormal vectors, hence a basis. Substitution into \eqref{eq:lap} gives
\[
 \AN v_{N,k}
 =h^{-2}\left(2-\ee^{-2\pi\ii k/N}-\ee^{2\pi\ii k/N}\right)v_{N,k}.
\]
Using \(2-2\cos(2q)=4\sin^2q\) proves \eqref{eq:spectrum}.
Only \(k=0\) gives zero in the selected range.
The identity \(\sin(\pi-\pi k/N)=\sin(\pi k/N)\) gives the pairing.
On \(0\leq k\leq N/2\) sine squared is strictly increasing, so there
are no further coincidences. At even \(N\) the midpoint pairs with itself.
\end{proof}

The equality \(v_{N,N-k}=\overline{v_{N,k}}\) explains the two directions.
The corresponding real eigenspace has the cosine and sine basis
\(\sqrt{2/L}\cos(2\pi kj/N)\),
\(\sqrt{2/L}\sin(2\pi kj/N)\). It has dimension two over \(\R\), just
as the two traveling-wave vectors have dimension two over \(\C\) in the
complexified space. The even midpoint instead has real entries
\(L^{-1/2}(-1)^j\) and multiplicity one.

\begin{figure}[ht]
\centering
\setlength{\unitlength}{1pt}
\begin{picture}(410,128)
\qbezier(133,60)(133,75.74)(121.87,86.87)
\qbezier(121.87,86.87)(110.74,98)(95,98)
\qbezier(95,98)(79.26,98)(68.13,86.87)
\qbezier(68.13,86.87)(57,75.74)(57,60)
\qbezier(57,60)(57,44.26)(68.13,33.13)
\qbezier(68.13,33.13)(79.26,22)(95,22)
\qbezier(95,22)(110.74,22)(121.87,33.13)
\qbezier(121.87,33.13)(133,44.26)(133,60)
\qbezier(348,60)(348,75.74)(336.87,86.87)
\qbezier(336.87,86.87)(325.74,98)(310,98)
\qbezier(310,98)(294.26,98)(283.13,86.87)
\qbezier(283.13,86.87)(272,75.74)(272,60)
\qbezier(272,60)(272,44.26)(283.13,33.13)
\qbezier(283.13,33.13)(294.26,22)(310,22)
\qbezier(310,22)(325.74,22)(336.87,33.13)
\qbezier(336.87,33.13)(348,44.26)(348,60)
\put(133,60){\makebox(0,0){\rule{3pt}{3pt}}}\put(122,87){\makebox(0,0){\rule{3pt}{3pt}}}\put(95,98){\makebox(0,0){\rule{3pt}{3pt}}}
\put(68,87){\makebox(0,0){\rule{3pt}{3pt}}}\put(57,60){\makebox(0,0){\rule{3pt}{3pt}}}\put(68,33){\makebox(0,0){\rule{3pt}{3pt}}}
\put(95,22){\makebox(0,0){\rule{3pt}{3pt}}}\put(122,33){\makebox(0,0){\rule{3pt}{3pt}}}
\put(348,60){\makebox(0,0){\rule{3pt}{3pt}}}\put(337,87){\makebox(0,0){\rule{3pt}{3pt}}}\put(310,98){\makebox(0,0){\rule{3pt}{3pt}}}
\put(283,87){\makebox(0,0){\rule{3pt}{3pt}}}\put(272,60){\makebox(0,0){\rule{3pt}{3pt}}}\put(283,33){\makebox(0,0){\rule{3pt}{3pt}}}
\put(310,22){\makebox(0,0){\rule{3pt}{3pt}}}\put(337,33){\makebox(0,0){\rule{3pt}{3pt}}}
\put(78,103){\makebox(34,12){\(\longleftarrow\)}}\put(293,103){\makebox(34,12){\(\longrightarrow\)}}
\put(55,4){\makebox(80,12){\(k=1\)}}
\put(261,4){\makebox(98,12){\(k=N-1\equiv-1\)}}
\put(62,53){\makebox(66,14){\(\ee^{\,2\pi\ii j/N}\)}}
\put(277,53){\makebox(66,14){\(\ee^{-2\pi\ii j/N}\)}}
\put(169,75){\makebox(67,14){conjugate modes}}
\put(169,55){\makebox(67,14){same eigenvalue}}
\put(177,33){\makebox(50,12){\(\longleftrightarrow\)}}
\end{picture}
\caption{Two opposite Fourier phase windings on the same discretized circle
(\(N=8\) drawn schematically). Arrows indicate the sign of phase winding
as the site index increases, not a postulated motion in time.
These are two modes, not two particles or two separate rings.}
\label{fig:pair}
\end{figure}

\section{The exact trace correspondence}\label{sec:trace}
\subsection{A convergent positive spectral observable}
For real \(p>0\) define
\begin{equation}
 Z_N(p;L)=\Tr_{\mathcal H_N^0}(\AN^{-p})
  =\sum_{k=1}^{N-1}\lambda_{N,k}(L)^{-p},
 \qquad T_N(L)=Z_N(3/2;L).
 \label{eq:trace}
\end{equation}
Functional calculus here simply means raising each positive eigenvalue
to \(-p\) in the Fourier basis. Although different inner-product
normalizations change the normalized vectors, they do not change these
eigenvalues or the trace.

\begin{theorem}[Paired-mode continuum limit]\label{thm:limit}
For every fixed \(L>0\) and real \(p>1/2\),
\begin{equation}
 \lim_{N\to\infty}Z_N(p;L)
  =2\left(\frac{L}{2\pi}\right)^{2p}\zeta(2p).
 \label{eq:general}
\end{equation}
In particular,
\begin{equation}
 \boxed{\displaystyle \lim_{N\to\infty}T_N(L)
             =\frac{L^3}{4\pi^3}\zeta(3).}
 \label{eq:limit}
\end{equation}
\end{theorem}
\begin{proof}
Set \(M_N=\lfloor(N-1)/2\rfloor\), and let \(\delta_N\) be one for even
\(N\) and zero for odd \(N\). The complete sum is
\begin{equation}
 Z_N(p;L)=2\sum_{k=1}^{M_N}\lambda_{N,k}(L)^{-p}
          +\delta_N\left(\frac{L}{2N}\right)^{2p}.
 \label{eq:paired}
\end{equation}
For \(1\leq k\leq N/2\), concavity of sine on \([0,\pi/2]\) places it
above its endpoint chord:
\(\sin(\pi k/N)\geq2k/N\). Consequently
\[
 \lambda_{N,k}(L)\geq16k^2/L^2,\qquad
 0\leq2\lambda_{N,k}(L)^{-p}
       \leq2(L/4)^{2p}k^{-2p}.
\]
Extend the paired summand by zero for \(k>M_N\). Its majorant is a
summable sequence independent of \(N\), because \(2p>1\).
For every fixed positive \(k\),
\[
 \lambda_{N,k}(L)
 =\left(\frac{2\pi k}{L}\right)^2
   \left(\frac{\sin(\pi k/N)}{\pi k/N}\right)^2
 \longrightarrow (2\pi k/L)^2.
\]
Dominated convergence for the series therefore gives twice the positive
integer sum of \((2\pi k/L)^{-2p}\).
The midpoint term in \eqref{eq:paired} is bounded by
\((L/2N)^{2p}\) and tends to zero. This proves \eqref{eq:general};
substitution \(p=3/2\) proves \eqref{eq:limit}.
\end{proof}

The pairing is necessary, not a cosmetic convention. A mode whose index is
\(N-k\), with \(k\) fixed, remains a low-frequency mode as \(N\) grows: on
the lattice it is exactly the mode with signed frequency \(-k\).
Looking only at each fixed positive index misses this second end of the
finite spectrum. Summing those one-sided limits yields half of
\eqref{eq:limit}. Equation~\eqref{eq:paired} handles the moving endpoint
before taking the limit.

The range \(p>1/2\) is also substantive. Since
\(\sin q\leq q\), the first \(M_N\) positive modes obey
\(\lambda_{N,k}^{-p}\geq(L/2\pi)^{2p}k^{-2p}\).
Thus the trace diverges for \(0<p\leq1/2\) as \(N\to\infty\).
No analytic continuation is involved in the positive convergent trace
used here.

\subsection{A quantitative finite-size control}
The convergence proof can be supplemented by a bound a reader can evaluate
at any resolution. Let \(\mathrm H_M=\sum_{k=1}^M k^{-1}\) be the harmonic
number, not a Hamiltonian.

\begin{proposition}[An explicit trace error bound]\label{prop:error}
With \(M=M_N\geq1\) and \(\delta_N\) as above,
\begin{equation}
 \left|T_N(L)-\frac{L^3\zeta(3)}{4\pi^3}\right|
 \leq L^3\left[
 \frac{\pi^2\mathrm H_M}{64N^2}
 +\frac{1}{8\pi^3M^2}+\frac{\delta_N}{8N^3}\right].
 \label{eq:error}
\end{equation}
In particular the error is \(O(L^3\log N/N^2)\).
\end{proposition}
\begin{proof}
For \(0<q\leq\pi/2\), integration of \(1-\cos q\leq q^2/2\)
gives \(0\leq q-\sin q\leq q^3/6\). Also \(\sin q\geq2q/\pi\).
It follows that
\begin{equation}
 0\leq\left(\frac q{\sin q}\right)^3-1
 =\frac{(q-\sin q)(q^2+q\sin q+\sin^2q)}{\sin^3q}
 \leq\frac{\pi^3q^2}{16}.
 \label{eq:sinbound}
\end{equation}
Using \(q=\pi k/N\), the excess of the paired finite sum over its
paired continuum truncation is therefore between zero and
\[
 \frac{L^3}{4\pi^3}
 \sum_{k=1}^M k^{-3}\frac{\pi^3}{16}
                         \left(\frac{\pi k}{N}\right)^2
 =\frac{L^3\pi^2}{64N^2}\mathrm H_M.
\]
The omitted continuum tail is at most
\((L^3/4\pi^3)\int_M^\infty x^{-3}dx=L^3/(8\pi^3M^2)\).
The remaining even midpoint is \(\delta_N L^3/(8N^3)\).
The triangle inequality proves \eqref{eq:error}, and
\(\mathrm H_M\leq1+\log M\) proves the stated order.
\end{proof}

This is a bound for the whole trace. It does not require uniform
convergence of every finite eigenvalue to a continuum eigenvalue.
Indeed, the most oscillatory modes remain far from the parabolic
continuum dispersion in relative terms. Their inverse-power contribution
is small enough for the complete trace to converge.

\subsection{Exact small cycles and finite illustrations}
The first three examples already test both parity cases:
\begin{align}
 T_3(L)&=\frac{2L^3}{81\sqrt3}, \label{eq:t3}\\
 T_4(L)&=\frac{L^3}{512}(1+4\sqrt2), \label{eq:t4}\\
 T_5(L)&=\frac{2L^3}{125}
 \left[\left(\frac{5-\sqrt5}{2}\right)^{-3/2}
      +\left(\frac{5+\sqrt5}{2}\right)^{-3/2}\right].
 \label{eq:t5}
\end{align}
For \(N=3\) the nonzero eigenvalue is \(27/L^2\) twice.
For \(N=4\) the eigenvalues are \(32/L^2,64/L^2,32/L^2\);
the midpoint \(64/L^2\) must be counted once.
For \(N=5\) each of \(25(5-\sqrt5)/(2L^2)\) and
\(25(5+\sqrt5)/(2L^2)\) occurs twice.
These statements follow directly from \eqref{eq:spectrum} and can also be
checked by characteristic polynomials.

Table~\ref{tab:finite} illustrates the approach to the limit at unit
circumference. All entries are evaluations of a fixed finite formula, not
evolving-state simulations. The proof is Theorem~\ref{thm:limit}.
\begin{table}[ht]
\centering
\begin{tabular}{rcc}
\toprule
\(N\)& \(T_N(1)\) & \(T_N(1)/[\zeta(3)/(4\pi^3)]\)\\
\midrule
3 & \(0.0142555622022\) & \(1.47085185\)\\
4 & \(0.013001668456\) & \(1.34147836\)\\
5 & \(0.0121735234086\) & \(1.2560325\)\\
6 & \(0.0116199082382\) & \(1.19891193\)\\
16 & \(0.0101164547892\) & \(1.04378951\)\\
32 & \(0.00982510692991\) & \(1.01372899\)\\
\(\infty\)& \(0.00969204490073\) & \(1\)\\
\bottomrule
\end{tabular}
\caption{Bounded finite-value illustrations, rounded from 60-digit
arithmetic. Values and rounding are reproduced by the supplied checker.
The selected rows do not assert monotonicity for arbitrary \(N\).}
\label{tab:finite}
\end{table}

\section{Correspondence with standard quantum ring theory}\label{sec:quantum}
\subsection{The continuum operator and its domain}
For a spinless free particle of mass \(m>0\), let
\(\mathcal H=L^2([0,L],dx)\). Take \(\hbar>0\) and the standard Hamiltonian
\begin{equation}
 H=-\frac{\hbar^2}{2m}\frac{d^2}{dx^2},\qquad
 D(H)=\{\psi\in H^2(0,L):
 \psi(0)=\psi(L),\ \psi'(0)=\psi'(L)\}.
 \label{eq:ham}
\end{equation}
Both periodic boundary conditions are part of the operator domain;
a value condition alone is not a full self-adjoint specification for a
second derivative. We consider the untwisted, zero-flux circle, with no
potential or internal spin multiplicity. The familiar ring calculation is
presented in standard teaching treatments \cite{zwiebach,mitring}.

The normalized waves
\begin{equation}
 \phi_k(x)=L^{-1/2}\ee^{2\pi\ii kx/L},\qquad k\in\Z,
 \label{eq:plane}
\end{equation}
satisfy both boundary conditions. Substitution into \eqref{eq:ham} gives
\begin{equation}
 E_k=\frac{\hbar^2}{2m}\left(\frac{2\pi k}{L}\right)^2.
 \label{eq:energy}
\end{equation}
Conversely, positivity of the integration-by-parts form excludes negative
eigenvalues. At zero the solution is affine and periodicity makes it
constant. At positive energy write the differential equation's solution
as \(a\ee^{\ii qx}+b\ee^{-\ii qx}\). Imposing both boundary conditions
allows a nonzero solution exactly when \(qL=2\pi n\), \(n\) a positive
integer. These are the two waves \(k=\pm n\).

Completeness is the Fourier-series theorem on the circle. More explicitly,
the Fourier transform identifies \(H\) with the diagonal real operator
\((c_k)\mapsto(E_kc_k)\) on
\(\{c:\sum_k E_k^2|c_k|^2<\infty\}\).
A diagonal real multiplication operator on its maximal domain is
self-adjoint; its domain here is precisely the periodic \(H^2\) domain.
The constant ground state has \(E_0=0\), while every \(E_k>0\) is doubly
degenerate. Momentum \(p_k=2\pi\hbar k/L\) has a sign, whereas the kinetic
energy depends on its square. The twofold degeneracy is directional,
not an additional particle count.

\begin{proposition}[Continuum inverse-energy trace]\label{prop:qtrace}
On the positive spectral subspace of \(H\), for real \(p>1/2\),
\begin{equation}
 \zeta_H(p):=\sum_{k\in\Z\setminus\{0\}}E_k^{-p}
 =\left(\frac{2m}{\hbar^2}\right)^p
     2\left(\frac{L}{2\pi}\right)^{2p}\zeta(2p).
 \label{eq:qtrace}
\end{equation}
\end{proposition}
\begin{proof}
Substitute \eqref{eq:energy}, pair \(k\) and \(-k\), and factor out
the constant. Absolute convergence follows from \(2p>1\).
\end{proof}
The trace is a property of the operator's spectrum; it is not a choice of
a prepared wavefunction or its detection probabilities.

\subsection{The discrete Hamiltonian and the exact finite dictionary}
We now make a declared adoption:
\begin{equation}
 H_N=\frac{\hbar^2}{2m}\AN.
 \label{eq:adopt}
\end{equation}
It has the same Fourier eigenvectors as \(\AN\), and exactly
\begin{equation}
 E_{N,k}=\frac{\hbar^2}{2m}\lambda_{N,k}(L),\qquad
 \zeta_{H_N}(p)=\left(\frac{2m}{\hbar^2}\right)^p Z_N(p;L).
 \label{eq:qfinite}
\end{equation}
The finite zeta sum excludes the single zero mode but includes every
nonzero multiplicity. Thus Theorem~\ref{thm:limit} immediately gives
\begin{corollary}[Quantum spectral correspondence]\label{cor:quantum}
For \(p>1/2\), \(\zeta_{H_N}(p)\to\zeta_H(p)\). In particular
\begin{equation}
 \boxed{\displaystyle
 \lim_{N\to\infty}\zeta_{H_N}(3/2)
 =\left(\frac{2m}{\hbar^2}\right)^{3/2}
   \frac{L^3\zeta(3)}{4\pi^3}.}
 \label{eq:qmain}
\end{equation}
\end{corollary}

This identifies precisely what is exact at finite \(N\) and what is a
limit. Equation~\eqref{eq:qfinite} is exact by the adoption
\eqref{eq:adopt}. The finite sine spectrum is not the continuum
parabolic spectrum. The signed lattice labels lie in a finite frequency
band and identify \(k\) modulo \(N\). In particular the even midpoint
cannot carry two independent opposite modes.

For a fixed nonzero integer \(k\) with \(|k|<N/2\), let its lattice index
be \(k\bmod N\). Its relative kinetic energy is
\begin{equation}
 \frac{E_{N,k\bmod N}}{E_k}
 =\left(\frac{\sin(\pi |k|/N)}{\pi |k|/N}\right)^2
 =1-\frac{\pi^2k^2}{3N^2}+O(k^4/N^4).
 \label{eq:dispersion}
\end{equation}
The expansion is for fixed \(k\), or for \(|k|/N\to0\) with the remainder
controlled accordingly. At the even midpoint the analogous ratio is
\(4/\pi^2\), independent of \(N\); refinement does not make the relative
error vanish for those moving high-frequency labels.

The units track the correspondence. If \(L\) is a length,
\([\lambda]=\mathrm{length}^{-2}\) and
\([T_N]=\mathrm{length}^3\).
Since \([\hbar]=\mathrm{mass}\,\mathrm{length}^2/\mathrm{time}\),
\([\hbar^2/(2m)]=\mathrm{energy}\,\mathrm{length}^2\).
Thus \(E_{N,k}\) has energy units and \(\zeta_{H_N}(3/2)\) has
units \(\mathrm{energy}^{-3/2}\). None of these traces is itself a
dimensionless probability.

\subsection{What the standard quantum time law adds}
Self-adjointness permits the conventional Schr\"odinger evolution
\begin{equation}
 U_N(t)=\exp(-\ii tH_N/\hbar),\qquad
 U_N(t)v_{N,k}=\ee^{-\ii tE_{N,k}/\hbar}v_{N,k}.
 \label{eq:unitary}
\end{equation}
Because every \(E_{N,k}\) is real, every spectral factor has modulus one.
Consequently \(U_N(t)^*U_N(t)=I\),
\(U_N(t+s)=U_N(t)U_N(s)\), and
\(\|U_N(t)u\|_h=\|u\|_h\).
Differentiating the finite exponential shows
\(\ii\hbar\,\partial_t U_N(t)u=H_NU_N(t)u\).
This is a proof of unitarity and of the time equation under its explicit
Hamiltonian adoption.

For \(u=\sum_k a_kv_{N,k}\), the coefficients evolve as
\(a_k(t)=a_k(0)\ee^{-\ii E_{N,k}t/\hbar}\).
The relative phase of two components changes by
\(-(E_{N,k}-E_{N,\ell})t/\hbar\).
The degenerate pair \(k,N-k\) acquires the same time phase, so its relative
phase is constant. A superposition of different energies can change its
spatial profile even though the norm is conserved.
These temporal statements follow from \eqref{eq:unitary}, not from
periodicity or a spatial trace alone.

There is also a controlled relation to continuum time evolution.
Let \(\mathcal B_K=\operatorname{span}\{\phi_k:|k|\leq K\}\), with \(K\)
a fixed nonnegative integer. If \(N>2K\), sampling at the sites defines an
isometry \(J_{N,K}:\mathcal B_K\to\C^N\),
\(J_{N,K}\phi_k=v_{N,k\bmod N}\), by Fourier orthogonality.

\begin{proposition}[Finite-band dynamical correspondence]\label{prop:band}
For \(N>2K\), \(t\in\R\), and \(\psi\in\mathcal B_K\),
\begin{equation}
 \|U_N(t)J_{N,K}\psi-J_{N,K}U(t)\psi\|_h
 \leq\frac{2\hbar\pi^4K^4|t|}{3mL^2N^2}\|\psi\|_{L^2},
 \qquad U(t)=\ee^{-\ii tH/\hbar}.
 \label{eq:band}
\end{equation}
\end{proposition}
\begin{proof}
For \(0\leq q\leq\pi/2\), \(\sin q\geq q-q^3/6\) implies
\(0\leq1-(\sin q/q)^2\leq q^2/3\), with the ratio defined by its
limit at zero. Equations~\eqref{eq:energy}--\eqref{eq:dispersion} give
\[
 0\leq E_k-E_{N,k\bmod N}
 \leq\frac{2\hbar^2\pi^4K^4}{3mL^2N^2}
 \quad (|k|\leq K).
\]
For real \(a,b\), integration of the derivative of \(\ee^{-\ii x}\)
gives \(|\ee^{-\ii a}-\ee^{-\ii b}|\leq|a-b|\).
Apply this to each time phase, square and sum in the orthonormal
sampled basis. Taking the square root proves \eqref{eq:band}.
\end{proof}

The estimate is uniform on a fixed frequency band and bounded time
interval. It is not a statement that all lattice modes approximate the
continuum uniformly, or that the historical SRG recursion evolves
unitarily. The latter would require an independently specified SRG
temporal operator with the necessary properties. Here \(t\), \(m\),
\(\hbar\), the energy interpretation and the Schr\"odinger law all enter
through the conventional quantum comparison.

\section{The historically selected SRG coupling}\label{sec:coupling}
The historical relation relevant here is
\begin{equation}
 \pi e\phig=\gamma^{-1}\zeta(3),\qquad
 \gamma_*=\frac{\zeta(3)}{\pi e\phig},\qquad
 \phig=\frac{1+\sqrt5}{2}.
 \label{eq:lock}
\end{equation}
It is printed in the July 2025 SRG sources
\cite[physical p.~3, Eq.~(1)]{p01}
\cite[physical p.~4, Eq.~(2)]{p06}.
We use the notation \(\gamma_*\) for the solved reference value, as in the
completed reconstruction. The historical record precedes the October
loop diagnostic: the circle trace did not cause the original printed
formula.

\begin{corollary}[Exact calibrated SRG identity]\label{cor:calibration}
With unit dimensionless circumference,
\begin{equation}
 \boxed{\displaystyle
 \gamma_*=\frac{4\pi^2}{e\phig}\lim_{N\to\infty}T_N(1).}
 \label{eq:calibration}
\end{equation}
For a dimensionful circumference the equivalent identity is
\begin{equation}
 \gamma_*=\frac{4\pi^2}{e\phig}
           \lim_{N\to\infty}\frac{T_N(L)}{L^3}.
 \label{eq:dimensionless}
\end{equation}
\end{corollary}
\begin{proof}
Insert \eqref{eq:limit}; the product
\((4\pi^2/(e\phig))(\zeta(3)/(4\pi^3))\)
is exactly \(\zeta(3)/(\pi e\phig)\).
The scaling \(T_N(L)=L^3T_N(1)\) follows directly from
\eqref{eq:spectrum}.
\end{proof}

The multiplier \(4\pi^2/(e\phig)\) is a calibration introduced to match
the historical constant. It is not independently selected by the
Laplacian or by a physical measurement. The equality is exact; the
normalization is not derived from first principles.
The loop contributes an intrinsic positive spectral number after its
circumference and exponent have been specified. The additional \(e\) and
\(\phig\) enter through the historical reference, not through new
eigenvalues of the circle.

If \(\ell_0\) is a physical reference length and \(\widehat L=L/\ell_0\),
then \(\widehat\lambda=\ell_0^2\lambda\) is dimensionless and its
inverse-power trace is \(T_N(L)/\ell_0^3\). Saying simply \(L=1\)
means either dimensionless coordinates from the outset or this explicit
choice of unit circumference. Dividing by \(L^3\) in
\eqref{eq:dimensionless} makes the comparison independent of the unit
chosen for length.

The quantum version can be written with no remaining mass dependence:
\begin{equation}
 \gamma_*=\frac{4\pi^2}{e\phig L^3}
 \left(\frac{\hbar^2}{2m}\right)^{3/2}
       \zeta_H(3/2).
 \label{eq:qcalibration}
\end{equation}
The factor involving \(m,\hbar\) exactly undoes the adopted kinetic scale.
Thus this identity supplies no selected particle mass; every \(m>0\)
satisfies it with its own standard ring Hamiltonian.
Its scientific content is the relation between a documented SRG
reference and a precisely identified spectral observable, with all
scales visible.

\section{Three spectral observables, three different questions}\label{sec:observables}
The Origins paper discusses a positive inverse trace, the inverse
\(3/2\) trace and a negative finite part \cite[section 5]{origins}.
They can be defined on the same spectrum without representing the same
measurement. This section records their distinctions through short
derivations.

\subsection{The inverse trace and the positive value \(1/12\)}
\begin{proposition}[Exact finite inverse trace]\label{prop:inverse}
For \(N\geq3\),
\begin{equation}
 Z_N(1;L)=\frac{L^2}{12}(1-N^{-2}).
 \label{eq:inverse}
\end{equation}
\end{proposition}
\begin{proof}
Factorization of \(z^N-1\), with \(z=\ee^{2\ii x}\), gives
\(\prod_{k=0}^{N-1}\sin(x+k\pi/N)=c_N\sin(Nx)\)
for a nonzero constant \(c_N\). Taking a logarithmic derivative away
from zeros, and differentiating once more, yields
\[
 \sum_{k=0}^{N-1}\csc^2(x+k\pi/N)=N^2\csc^2(Nx).
\]
Subtract \(\csc^2x\) and let \(x\to0\).
The expansion \(\csc^2x=x^{-2}+1/3+O(x^2)\) leaves
\(\sum_{k=1}^{N-1}\csc^2(\pi k/N)=(N^2-1)/3\).
Multiplication by \(L^2/(4N^2)\) proves \eqref{eq:inverse}.
\end{proof}

The dimensionless limit \(1/12\) is an inverse-square-frequency sum.
By contrast \(T_N(L)/L^3\) is the inverse-cube-frequency sum giving
\(\zeta(3)/(4\pi^3)\). Changing the exponent changes the observable
even while leaving the geometry untouched.

\subsection{The negative finite part}
Define the dimensionless half-total frequency statistic
\begin{equation}
 S_N=\frac{L}{4\pi}\sum_{k=1}^{N-1}\sqrt{\lambda_{N,k}(L)}
     =\frac N{2\pi}\sum_{k=1}^{N-1}\sin(\pi k/N).
 \label{eq:freq}
\end{equation}
The square roots here are wave-number magnitudes. They become physical
oscillation frequencies only after an additional dispersion or time law
has been specified. In particular the free quantum particle has time
frequencies \(E_k/\hbar\), quadratic rather than linear in \(|k|\).

\begin{proposition}[A specified bulk subtraction]\label{prop:finitepart}
With \(B_N=N^2/\pi^2\),
\begin{equation}
 S_N=\frac N{2\pi}\cot\frac{\pi}{2N},\qquad
 S_N-B_N=-\frac1{12}-\frac{\pi^2}{720N^2}+O(N^{-4}).
 \label{eq:finitepart}
\end{equation}
\end{proposition}
\begin{proof}
The imaginary part of a finite geometric sum gives
\[
 \sum_{k=1}^{N-1}\sin(kq)
 =\frac{\sin((N-1)q/2)\sin(Nq/2)}{\sin(q/2)}.
\]
Set \(q=\pi/N\) to obtain the cotangent.
The continuous bulk for the same sine dispersion is
\[
 \frac N{2\pi}\int_0^N\sin(\pi x/N)\,dx=N^2/\pi^2.
\]
Finally \(\cot y=y^{-1}-y/3-y^3/45+O(y^5)\) gives
\eqref{eq:finitepart}.
\end{proof}

Every summand in \(S_N\) is positive. Its negative remainder has a simple
geometric sign: the composite trapezoidal sum for the strictly concave
sine arch lies below the integral, and both endpoint values are zero.
Thus \(S_N<B_N\) for every \(N\geq3\).
The limit of this difference equals the analytically continued value
\(\zeta(-1)=-1/12\), as recorded by the standard integer-value formula
for the zeta function \cite[Eq.~25.6.3]{dlmf}.
The ordinary sum \(\sum_{n\geq1}n\) still diverges to \(+\infty\).

This distinction also has an operator interpretation. The positive
inverse traces are ordinary convergent sums on the mean-zero space.
The statistic in \eqref{eq:freq} grows with the high-frequency modes and
requires its stated bulk subtraction to have the displayed limit.
It is not the inverse-energy trace of the quantum particle, nor an
energy expectation in a normalized state.
In particular no vacuum energy, oscillator zero-point factor or
regularization prescription is silently added to the SRG coupling.

\subsection{A heat-trace reading of the positive moment}
A useful connection to spectral geometry is obtained without any
subtraction beyond removing the zero mode. For \(\sigma>0\), let
\[
 K_N^0(\sigma)=\sum_{k=1}^{N-1}\ee^{-\sigma\lambda_{N,k}(L)}.
\]
The elementary gamma integral, after \(y=\sigma\lambda\), gives
\begin{equation}
 Z_N(p;L)=\frac1{\Gamma(p)}
     \int_0^\infty \sigma^{p-1}K_N^0(\sigma)\,d\sigma,\qquad p>0.
 \label{eq:mellin}
\end{equation}
The finite sum can be integrated term by term.
For the continuum circle, the identical statement holds with
\(K^0(\sigma)=2\sum_{k\geq1}\exp[-\sigma(2\pi k/L)^2]\)
when \(p>1/2\): nonnegativity permits interchange by Tonelli's theorem,
and the resulting sum is finite by Theorem~\ref{thm:limit}'s
continuum expression. This is the standard heat-trace interpretation
of a spectral zeta function \cite[section 2.1]{karlsson}.

Here \(\sigma\) has units of length squared. It is an integration
parameter, not a recovered SRG clock or the real time \(t\) of
\eqref{eq:unitary}. The Mellin form explains why an inverse moment
aggregates response across geometric scales. It also makes the
zero-mode removal transparent: retaining the constant term would make
the large-\(\sigma\) integral divergent.

\section{Provenance and status of the construction}\label{sec:provenance}
This paper uses the contemporary expansion \emph{Symbolic Resonance
Geometry}, consistent with SRG-I, while retaining actual historical
titles in the bibliography. The source corpus contains more than one
operator family and more than one interpretation of its symbols
\cite{srgi}. Our construction takes the normalized loop from the
October reconstruction; it does not replace any of those sources.

The inspected editions are \emph{SRG Origins} v0.1, dated 11 October
2026; \emph{SRG-I: Foundational Reconstruction} v0.1.1, dated
10 October 2026; and \emph{SRG Multiplication Geometry: Factor Data,
Circle Construction, and Mass-Shell Correspondence}, v0.1, dated 13 October 2026
\cite{origins,srgi,massshell}.
The last provides methodological context for a conditional physical
dictionary; its mass-shell geometry is not used to infer the ring
Hamiltonian or its mass.

The July 15 fusion source calls the lock empirically calibrated and
explicitly says it was not obtained from first principles
\cite[physical pp.~3--4]{p01}. This is a recovered historical claim,
not new measurement evidence supplied here.
The July 13 source prints the same lock but also calls its gamma the
Euler--Mascheroni constant in its glossary
\cite[physical pp.~4,6]{p06}. The solved value in \eqref{eq:lock},
approximately \(0.0869947475\), is different from that constant.
We retain the source discrepancy and use \(\gamma_*\) only for the
solved reference.

\begin{table}[ht]
\small
\centering
\begin{tabularx}{\textwidth}{@{}p{11mm}>{\raggedright\arraybackslash}p{40mm}>{\raggedright\arraybackslash}X@{}}
\toprule
Status & Item & Role in this paper\\
\midrule
H & Historical lock and source terminology &
July 2025 printed definitions and attributed calibration claims;
not a result of the October loop construction.\\
A & Normalized loop diagnostic &
Fixed \(L\), refinement by \(N\), positive sign, zero-mode exclusion
and exponent \(3/2\), adopted in the October reconstruction.\\
D & Spectrum and trace limits &
Consequences of that declared diagnostic; rederived here with
multiplicity, finite-size bounds and units explicit.\\
A & Historical calibration &
The factor \(4\pi^2/(e\varphi)\) matches the preexisting reference.\\
A & Quantum Hamiltonian and time law &
The conventional choices \(H_N=(\hbar^2/2m)\AN\) and
\(U_N(t)=\exp(-\ii tH_N/\hbar)\).\\
D & Quantum correspondence &
Exact energy scaling, trace convergence and finite-band estimate
after the declared adoption.\\
U & Native physical selection &
An independent SRG statistic selecting the calibration,
physical mass, clock and measurement correspondence remains
unspecified by these inputs.\\
\bottomrule
\end{tabularx}
\caption{H: recovered history; D: deduction; A: later diagnostic or
adoption; U: unresolved source or physical identification.
An adoption can have exact mathematical consequences.}
\label{tab:status}
\end{table}

The inspected repository HEAD is
\texttt{5f08666798f1c7a6b6a60bbf681d7367088f7e7a}.
Source identities are bound to the actual local editions by the
accompanying source map and SHA-256 records, rather than inferred from
HEAD alone. Existing tracked UI/workflow differences were recorded
before drafting; the manuscript does not use them as scientific inputs.
Prior review or publication status is not transferred to this new draft.

The main inverse-trace theorem and calibrated identity restate completed
Origins mathematics in a self-contained form. The explicit error bound,
finite-band time estimate and organized kinetic correspondence are
supporting derivations of this edition. They are applications of standard
spectral tools, without a claim of priority over the broader literature.
No Tri-Octagon or QI01--QI03 recurrence, synchronizer or result is an
input to the paper.

\section{Scientific significance and physical interpretation}\label{sec:discussion}
The correspondence identifies a common mathematical structure at three
levels. First, periodicity organizes both the discrete and continuous
state spaces into signed Fourier modes. Second, local differences and
spatial differentiation assign compatible low-frequency quadratic
responses. Third, a summable inverse spectral moment carries this
compatibility to an exact scalar limit. The last point is stronger than
a resemblance between two formulas: it includes all modes, their
multiplicity, a zero-mode convention and a controlled limiting argument.

From the SRG side, this gives a precise spectral meaning to the zeta
factor already present in the historical reference. A finite periodic
diagnostic can approach that reference after its stated calibration.
From the quantum side, it is the positive inverse-energy moment of the
standard free ring Hamiltonian after a transparent change of units.
The geometric result can be studied on its own, and the quantum
comparison can be tested mathematically without inserting a quantum
time law into the SRG historical record.

Several features make the comparison useful for modeling. The circumference
is a measurable geometric scale in the conventional ring model; changing it
rescales energies by \(L^{-2}\) and the geometric inverse \(3/2\) trace
by \(L^3\). The finite lattice exposes where a continuum approximation
is accurate and where its dispersion fails. The two directions explain
degeneracy without adding degrees of freedom by hand, and the
finite-band estimate quantifies how energy error accumulates as phase
error over time. These are usable structural constraints for any later
model that deliberately adopts this diagnostic.

A physical realization needs further identifications. One must specify
what physical system is confined to the ring, how its mass and
circumference are obtained, why potential and flux terms can be
neglected, and how the chosen states and observables relate to
preparation and measurement. The Born interpretation belongs to
standard quantum theory when that theory is adopted; it is not
deduced by computing a positive spectral trace.
The trace itself is neither a normalized state nor a detector rule.

This single-particle construction does not establish entanglement,
Bell correlations, computational advantage or superluminal
communication. Those claims require different state spaces, observables
and evidence. Keeping the present correspondence at its actual strength
preserves SRG as an independent mathematical research program: its
historical definitions remain intact, while a well-defined modern
diagnostic is placed in contact with a substantial body of established
physics.

\section{Conclusion}\label{sec:conclusion}
The normalized periodic-loop diagnostic used in the SRG reconstruction
has an exact inverse spectral-moment limit,
\(L^3\zeta(3)/(4\pi^3)\), with both Fourier directions included.
It is the same moment of the positive continuum circle Laplacian.
Adopting the conventional kinetic factor \(\hbar^2/(2m)\) turns this
geometric identity into a precise correspondence of discrete and
continuum quantum-ring spectral zeta functions. The finite spectra,
their low-frequency error and the resulting standard unitary phases
have all been derived explicitly.

The historical SRG reference is an exactly calibrated multiple of this
limit. Its prefactor remains a documented matching choice rather than
a new spectral selection principle. The inverse traces and the
negative finite part answer different questions about the same
spectrum, and their signs and dimensions follow from their definitions.

The result supplies a rigorous point of contact between SRG's
reconstructed spectral geometry and quantum ring mathematics.
A subsequent physical model could deliberately use the derived
Hamiltonian dictionary and its finite-size controls; doing so would be
a separately specified modeling step. The present paper completes the
mathematical correspondence and leaves the historical research unchanged.

\appendix
\section{Verification and reproducibility}\label{app:verification}
The review bundle contains the complete LaTeX source, compiled PDF,
a reading guide, a source/provenance map, a theorem-to-proof ledger,
an isolated exact checker and machine-readable results.
The checker imports symbolic mathematics and evaluates only declared
finite formulas. It imports no historical program or production kernel
and runs no trajectory or parameter search.

Exact checks cover representative characteristic polynomials,
Fourier eigenvalue algebra, odd/even multiplicity, the closed small-cycle
traces, scale transformations, the elementary inequalities' algebraic
steps, the low-frequency expansion, the quantum scale factors,
the finite-band coefficient, the inverse-trace identity and the
finite-part expansion. A separate category records 60-digit finite
evaluations for the fixed values \(N=3,4,5,6,16,32\).
The proof of an infinite limit is the summable-majorant argument in
Section~\ref{sec:trace}, not an extrapolation from those numbers.

The artifact QA record covers compilation, equation and citation
resolution, extracted text, font embedding and visual inspection of
the rendered pages. Source-state records compare the inspected editions,
historical originals and production files before and after preparation.
The evidence distinguishes exact identities, numerical illustrations
and document checks instead of treating them as interchangeable
validation.

\subsection*{Assistance and review status}
Hilmir Fr\'imann Halld\'orsson is the human author.
GPT supplied the scientific work order and is the designated reviewer.
Codex assisted with source inspection, proof development, literature
verification, manuscript preparation and artifact QA.
These assistance roles do not constitute coauthorship, external peer
review or institutional endorsement. This is a review draft; no publication
or release has been performed. Existing source-specific rights remain
with their sources; no new blanket license is assigned here.

\begin{thebibliography}{10}
\bibitem{origins}
H.~F.~Halld\'orsson,
\emph{SRG Origins---From Geometry to Recursive Operators},
v0.1, 11 October 2026, sections 2, 5 and 6; local preserved edition.
\href{https://github.com/pzychozen/trioctagon-physics/tree/main/SRG_REDISCOVERY/PAPERS/SRG_ORIGINS_FROM_GEOMETRY_TO_RECURSION/publication/v0.1}{Repository edition}.
Exact inspected identities are in the accompanying source map.
\bibitem{srgi}
H.~F.~Halld\'orsson,
\emph{SRG-I: Foundational Reconstruction}, v0.1.1,
10 October 2026; source index and historical-definition discussion.
\href{https://github.com/pzychozen/trioctagon-physics/tree/main/SRG_REDISCOVERY/PAPERS/SRG_I_FOUNDATIONAL_RECONSTRUCTION/publication/v0.1.1}{Repository edition}.
\bibitem{massshell}
H.~F.~Halld\'orsson,
\emph{SRG Multiplication Geometry: Factor Data, Circle Construction,
and Mass-Shell Correspondence},
v0.1, 13 October 2026, section 8 (conditional physical dictionary).
\href{https://github.com/pzychozen/trioctagon-physics/tree/main/SRG_REDISCOVERY/PAPERS/SRG_MULTIPLICATION_MASS_SHELL/publication/v0.1}{Repository edition}.
\bibitem{p01}
H.~F.~Halld\'orsson,
\emph{Symbolic Recursion and SRG Fusion: Emergent Lattice Fields from
Glyph Memory Dynamics}, 15 July 2025; physical pp.~3--4,
Eq.~(1), and pp.~19--20.
\href{https://zenodo.org/records/17253999}{Historical source deposit},
file \texttt{1.Symbolic Recursion and SRG Fusion.pdf}.
\bibitem{p06}
H.~F.~Halld\'orsson,
\emph{Recursive Dark Energy and Vacuum Feedback: A Vesica-Phase
Cosmological Framework}, 13 July 2025; physical p.~4, Eq.~(2),
and p.~6 glossary.
\href{https://zenodo.org/records/15872156}{Historical source deposit},
file \texttt{dark\_energy.pdf}.
\bibitem{friedli}
F.~Friedli and A.~Karlsson,
\emph{Spectral zeta functions of graphs and the Riemann zeta function
in the critical strip}, arXiv:1410.8010v2 (2015),
pp.~2--3, Eqs.~(1)--(2), and section 2.
\href{https://arxiv.org/abs/1410.8010v2}{Verified preprint edition}.
\bibitem{karlsson}
A.~Karlsson, \emph{Spectral zeta functions},
arXiv:1907.01832v1 (2019), section 2.1.
\href{https://arxiv.org/abs/1907.01832v1}{Verified preprint edition}.
\bibitem{zwiebach}
B.~Zwiebach, \emph{Quantum Physics I}, MIT OpenCourseWare,
8.04, Spring 2016, Lecture Note 10, 14 March 2016, section 3,
physical pp.~6--8 (and section 1 for time phases):
Solving the Time-Independent Schr\"odinger Equation;
Free Particle on a Circle.
\href{https://ocw.mit.edu/courses/8-04-quantum-physics-i-spring-2016/resources/mit8_04s16_lecnotes10/}{Official lecture notes}.
\bibitem{mitring}
MIT Physics 8.06, \emph{Midterm 2007 solutions},
problem 4(a), physical p.~5: Particle on a Ring.
\href{https://stuff.mit.edu/afs/athena/course/8/8.06/spring08/handouts/midterm2007_sol.pdf}{Official solution PDF}.
The angular-coordinate normalization is converted here to arc length.
\bibitem{dlmf}
NIST Digital Library of Mathematical Functions,
section 25.6, Integer Arguments, Eq.~25.6.3.
\href{https://dlmf.nist.gov/25.6.E3}{Official formula}.
External references checked 13 October 2026.
\end{thebibliography}
\end{document}
